% ===== Джерела до розділу 3. Ключі з префіксом ch03: =====
% Позиції ядра (bib-core.tex) тут НЕ дублюються.

% --- Теоретико-обчислювальні межі ---
\bibitem{ch03:cohen} Cohen F. Computer Viruses: Theory and Experiments. \emph{Computers \& Security}. 1987. Vol.~6, No.~1. P.~22--35. DOI: 10.1016/0167-4048(87)90122-2.

\bibitem{ch03:chess} Chess D. M., White S. R. An Undetectable Computer Virus. \emph{Proceedings of the Virus Bulletin Conference}. 2000.

\bibitem{ch03:spinellis} Spinellis D. Reliable Identification of Bounded-length Viruses is NP-complete. \emph{IEEE Transactions on Information Theory}. 2003. Vol.~49, No.~1. P.~280--284. DOI: 10.1109/TIT.2002.806137.

\bibitem{ch03:rice} Rice H. G. Classes of Recursively Enumerable Sets and Their Decision Problems. \emph{Transactions of the American Mathematical Society}. 1953. Vol.~74, No.~2. P.~358--366. DOI: 10.1090/S0002-9947-1953-0053041-6.

\bibitem{ch03:selcuk} Selçuk A. A., Orhan F., Batur B. Undecidable Problems in Malware Analysis. \emph{2017 12th International Conference for Internet Technology and Secured Transactions (ICITST)}. Cambridge: IEEE, 2017. P.~494--497.

% --- Статистичне виявлення ---
\bibitem{ch03:kay} Kay S. M. Fundamentals of Statistical Signal Processing. Vol.~II: Detection Theory. Upper Saddle River: Prentice Hall, 1998. 560~p.

\bibitem{ch03:fawcett} Fawcett T. An Introduction to ROC Analysis. \emph{Pattern Recognition Letters}. 2006. Vol.~27, No.~8. P.~861--874. DOI: 10.1016/j.patrec.2005.10.010.

\bibitem{ch03:axelsson} Axelsson S. The Base-Rate Fallacy and the Difficulty of Intrusion Detection. \emph{ACM Transactions on Information and System Security}. 2000. Vol.~3, No.~3. P.~186--205. DOI: 10.1145/357830.357849.

% --- Теорія інформації та подібність ---
\bibitem{ch03:shannon} Shannon C. E. A Mathematical Theory of Communication. \emph{Bell System Technical Journal}. 1948. Vol.~27, No.~3. P.~379--423.

\bibitem{ch03:lyda} Lyda R., Hamrock J. Using Entropy Analysis to Find Encrypted and Packed Malware. \emph{IEEE Security \& Privacy}. 2007. Vol.~5, No.~2. P.~40--45. DOI: 10.1109/MSP.2007.48.

\bibitem{ch03:kornblum} Kornblum J. Identifying Almost Identical Files Using Context Triggered Piecewise Hashing. \emph{Digital Investigation}. 2006. Vol.~3 (Supplement). P.~91--97. DOI: 10.1016/j.diin.2006.06.015.

\bibitem{ch03:kolter} Kolter J. Z., Maloof M. A. Learning to Detect and Classify Malicious Executables in the Wild. \emph{Journal of Machine Learning Research}. 2006. Vol.~7. P.~2721--2744.

\bibitem{ch03:li} Li M., Vitányi P. An Introduction to Kolmogorov Complexity and Its Applications. 3rd ed. New York: Springer, 2008. 790~p.

% --- Українські джерела ---
\bibitem{ch03:opirskyi-ml} Опірський І. Р., Мельничук М. Р. Застосування машинного навчання в статичному аналізі безпеки коду. \emph{Кібербезпека: освіта, наука, техніка}. 2026. URL: \url{https://csecurity.kubg.edu.ua/index.php/journal/article/view/1193}.

\bibitem{ch03:hybrid} Метод захисту програмного забезпечення на основі гібридного аналізу коду. \emph{Кібербезпека: освіта, наука, техніка}. 2025. URL: \url{https://csecurity.kubg.edu.ua/index.php/journal/article/view/871}.
